\textbf{Proof.}
Because both branches of the regression fit are untruncated and linear in the pseudo-outcomes, substitution of the displayed smoother coefficients gives the exact algebraic identity
\[
\widehat\tau^{\mathrm{sp}}_{h,J_{\mu,n},J_{r,n}}=\sum_{i=1}^n\Gamma_{i,h}Z_i.
\]
Write \(\epsilon_i=Z_i-\mu(X_i)\).  Add and subtract \(\sum_i\Gamma_{i,h}\mu(X_i)\) and \(m^{-1}\sum_j\mu(X_j^T)\).  The resulting stochastic part is exactly
\[
\sum_i\Gamma_{i,h}\epsilon_i+\frac1m\sum_j\{\mu(X_j^T)-\tau\}.
\]
The deterministic conditional-mean remainder splits into the boundary approximation error, the product of the interior regression and ratio errors, and the two ridge perturbations.  Indeed, on \([h,1]\), the identity \(f_T=x^{-a}qf_S\) gives
\[
\int_h^1(\widehat\mu-\mu)f_T
-\int_h^1\widehat r^{\mathrm{DS}}(\widehat\mu-\mu)f_S
=\int_h^1(\widehat q^{\mathrm{DS}}-q)(\mu-\widehat\mu)g_S.
\]
Cauchy--Schwarz and the two localized nuisance inequalities thus bound its root mean square by
\[
C\{\delta_{s,h}+\lambda_nJ_{\mu,n}\}
  \{\delta_{t,h}+\lambda_nJ_{r,n}\}.
\]
Since \(\lambda_n=n^{-2}\) and both dimensions are at most \(n\), the cross-products containing a ridge factor are bounded by \(C\lambda_n(J_{\mu,n}+J_{r,n})\).  On \([0,h]\), the order-\(p\) Taylor remainder for each of \(\mu_0\) and \(\mu_1\), integrated against \(f_T\le c_{T,+}\), is at most \(2C_{\mathrm{LP}}Mc_{T,+}h^{s+1}\).  These calculations prove the remainder bound.

For the coefficients, local support, the good-Gram bounds, and the partition of unity show that the regression-induced correction to the direct coefficient has the same quadratic order and smaller maximum order than the direct term.  The boundary smoother contributes
\[
\sum_i(\Gamma^{\mathrm{bd}}_{i,h})^2\asymp_P n^{-1}h^{1-a},
\]
while the interior term contributes
\[
\sum_i(\Gamma^{\mathrm{int}}_{i,h})^2\asymp_P n^{-1}\int_h^1x^{-a}\,dx.
\]
Their overlap is confined to \([h,2h]\); the Gram norm equivalence and Cauchy--Schwarz bound its cross term by the geometric mean, so the sum is comparable to \(n^{-1}V_a(h)\).  The corresponding maximum-to-sum ratio is bounded by the leverage expression proved in \(\mathrm{lem:localized\text{-}direct\text{-}series\text{-}control}\).  Taking \(h=h_n\) proves the last two assertions. \(\square\)